\chapter{Einleitung}
\pagenumbering{arabic}

%Platzhalter Einleitung

\section{Problemstellung}

\textbf{Überarbeiten wenn überhaupt verwenden ?!}

Konventionelle Membranlautsprecher weisen inhärente Limitationen auf: gerichtete Schallabstrahlung mit begrenztem Sweet Spot sowie voluminöse Gehäuse mit Bautiefen über 20 cm limitieren Designfreiheit und Raumintegration.
DML - Distributed-Mode-Loudspeaker bieten als Alternative omnidirektionale Abstrahlung (150-170°) und flache Bauformen (<2 cm) durch Biegewellen-Anregung in Platten \cite{Harris2000}. Dies eröffnet Anwendungen in Wandintegration, Automotive-Audio und architektonischen Installationen. Die zentrale Herausforderung besteht jedoch in Frequenzgang-Unregelmäßigkeiten von ±8-12 dB \cite{Bai2001}, welche aus komplexer modaler Struktur resultieren. Die akustische Performance wird maßgeblich durch die Exciter-Position determiniert.
Trotz existierender Hersteller-Empfehlungen fehlt eine systematische, Finite-Elemente-Methode (FEM)-gestützte Optimierung für spezifische Panel-Materialien. Zudem besteht Forschungsbedarf bezüglich der optimalen Systemarchitektur unter Berücksichtigung der Vokalbereichswiedergabe.
Für die Ostfalia Hochschule ergibt sich die Relevanz, durch Integration alternativer Schallwandler-Konzepte das Lehrangebot zu erweitern und die Innovationskompetenz angehender Ingenieure zu fördern. Es besteht somit der Bedarf an systematischer DML-Optimierung mittels Simulation und Experiment sowie der Entwicklung eines kostengünstigen, reproduzierbaren Prototyps für Lehrzwecke.


\section{Zielsetzung und Vorgehensweise}

\textbf{Überarbeiten}

Die vorliegende Arbeit verfolgt sowohl technische als auch pädagogische Zielsetzungen. Im technischen Kontext gilt es, einen funktionsfähigen Flachlautsprecher auf Basis der DML-Technologie zu entwickeln, welcher den Anforderungen der DIN 45500 für HiFi-Lautsprecher entspricht. Parallel dazu soll die Arbeit zur Erweiterung des Lehrangebots der Ostfalia Hochschule für angewandte Wissenschaften beitragen, indem Studierenden eine erweiterte Varietät von Akustik-Geräten geboten wird, um das technische Verständnis zu schärfen sowie Innovationskraft und Pioniergeist der zukünftigen Ingenieurgeneration zu fördern.

Im Zentrum der technischen Zielsetzung steht die systematische Entwicklung der Exciter-Positionierung auf einem Flachpanel sowie die optimale Größe der Panele bei gegebenen Materialien. Die Positionierung der Körperschallwandler aber auch die Größe und Eigenschaften der Panele beeinflusst maßgeblich die Anregung der Schwingungsmoden und beeinflusst folglich den resultierenden Frequenzgang. Diese Arbeit adressiert die bestehende Forschungslücke durch einen integrierten methodischen Ansatz, welcher experimentelle Freifeld-Messungen, Finite-Elemente-Simulationen in COMSOL Multiphysics sowie den Bau eines funktionsfähigen Flachlautsprechers kombiniert.

Übergeordnet leistet diese Arbeit einen Beitrag zur Erweiterung des an der Ostfalia verfügbaren Spektrums von Lautsprechertechnologien und bietet angehenden Ingenieuren einen fundierten Einblick in alternative Schallwandler-Konzepte. Durch die Auseinandersetzung mit innovativen Technologien wird die Innovationskompetenz und das kritische analytische Denken der Studierenden gefördert, was für ihre zukünftige Tätigkeit als Ingenieure von fundamentaler Bedeutung ist.

